The Catalan constant often appears as an alternating series:\par

\[
G
=
1-\frac1{3^2}+\frac1{5^2}-\frac1{7^2}+\cdots.
\]

That formula gives its value, and the HMT genealogy explains why this specific alternation exists.\par

In HMT, the origin is in the torsion chamber.\par

There exists an internal operator \(J\) such that\par

\[
J^2=-I.
\]

The root of \(-1\) emerges from two successive applications of a
quarter turn. The complex structure is the algebraic memory of an internal
rotation.\par

That turn generates a character of order four. When it is published on the positive integers, it produces exactly the residual pattern that alternates among \(+1\), \(0\), \(-1\), \(0\).\par

The Catalan constant then appears as a spectral trace:\par

\[
G
=
L(2,\chi_{-4})
=
\operatorname{Tr}(N^{-2}Q_4).
\]

This changes its meaning. Catalan's constant is simultaneously the sum of a particular series and the spectral response of the integers to a quarter of torsion.\par

Each term of the series records how the internal rotation acts on a residue class. The alternation is the arithmetic shadow of the rotation.\par

The triadic structure has its own character modulo \(3\). Combining it with the character of order four, the Chinese remainder theorem produces a character modulo \(12\):\par

\[
\chi_{12}
=
\chi_{-3}\chi_{-4}.
\]

The module \(12\) appears as the minimal common space in which triadic orientation and the quarter-torsion can coexist.\par

The dodecaphase is, in that sense, a reconciliation of \(3\) and \(4\).\par

Of the combined character emerges\par

\[
L(1,\chi_{12})
=
\frac{\log(2+\sqrt3)}{\sqrt3},
\]

y\par

\[
L(2,\chi_{12})
=
\frac{\pi^2}{6\sqrt3}.
\]

La cantidad\par

\[
2+\sqrt3
\]

is the fundamental unit of the real field \(\mathbb Q(\sqrt3)\). Its logarithm represents a hyperbolic shift. Once again, a finite residue structure opens an infinite multiplicative dynamics.\par

Here rotation, torsion, arithmetic character, Pell unit, and hyperbolic displacement meet.\par

The Euler–Kronecker constant of the field \(\mathbb Q(\sqrt3)\) prolongs this structure. It measures the regular term that remains after separating the principal pole of the field's zeta function. In the HMT genealogy, this term constitutes the regular invariant of fine arithmetic memory that remains after removing the universal divergence.\par

The Bendersky–Glaisher family extends the mechanism into a tower. The spectral derivatives of zeta generate that tower. Apery's constant and other constants associated with odd values are organized as members of the same functional family.\par

The novel contribution lies in the genealogy that unifies the classical constants:\par

\begin{itemize}[leftmargin=*,itemsep=0.35em]
\item the quarter of torsion produces the complex structure;
\item the residue triadico produces the structure modulo \(3\);
\item their composition produces the character modulo \(12\);
\item the character generates values \(L\);
\item the values \(L\) open the quadratic field and its spectral constants.
\end{itemize}

Geometric torsion has become arithmetic while preserving entirely its orientation.\par
